\section{Mathematical modelling}\label{sec:math}
\subsection{Causal information and forecast targets}
Table~\ref{tab:notation} defines the principal symbols and units. The model distinguishes physical validity from channel availability; the former controls scoring and the latter controls service access.

Let $G=(V,E)$ represent the sensor graph, $y_{i,u}$ the speed at sensor $i$ and target bin $u$, and $t$ an issue bin. The reporting channel assigns input arrival $a^x_{i,u}$ and feedback arrival $a^f_{i,u}$. Permanent loss is represented by an infinite arrival time. A value may enter the predictor only when $a^x_{i,u}\le t$ and may update calibration only when $a^f_{i,u}\le t$. The service information set is
\begin{equation}
\mathcal I_t=\{x_{i,u}:a^x_{i,u}\le t\}\cup\{y_{i,u}:a^f_{i,u}\le t\}\cup\{\text{previously issued records}\}.
\end{equation}
\label{eq:information}
The predictor $f_\theta$ is frozen and depends only on $\mathcal I_t$. For horizon $h$, define its centre $m_{i,t,h}$ and scale
\begin{equation}
m_{i,t,h}=\widehat q_{0.50,i,t,h},\qquad
s_{i,t,h}=\max\{(\widehat q_{0.95,i,t,h}-\widehat q_{0.05,i,t,h})/2,s_0\},
\label{eq:scale}
\end{equation}
where $s_0=0.44704\,\ms$, equivalent to the implementation's one-mile-per-hour floor. All reported dimensional scores are converted to SI units using the exact factor $1\,\mathrm{mph}=0.44704\,\ms$; this rescaling does not change coverage or comparison signs.

\subsection{Residual eligibility and relevance weighting}
A signed residual becomes eligible only after its target report arrives:
\begin{equation}
r_{i,t,h}=(y_{i,t+h}-m_{i,t,h})/s_{i,t,h},\qquad a^f_{i,t+h}\le t',
\end{equation}
where $t'$ is the current calibration time. The original issue-time centre and scale are used. For each horizon, the live archive $R_h(t')$ retains received residuals whose targets are at least $t'-B$, with $B=2016$ five-minute bins. Target age, rather than delivery age, governs this eligibility.

Write $c_{i,t}$ for a two-dimensional issue context: neighbourhood mean availability and the logarithm of one plus neighbourhood mean observation age, each standardised using training histories. For a pool $P$, a record $k$ with target $u_k$ receives unnormalised weight
\begin{equation}
w_k(t',c)=\exp\left[-\frac{t'-u_k}{\tau}-\frac{\|c_k-c\|_2^2}{2b^2}\right],\qquad
\widetilde w_k=w_k/\sum_{j\in P}w_j.
\label{eq:weights}
\end{equation}
The locked settings are $\tau=288$ bins and $b=2$. Grouping weights by target timestamp avoids counting correlated spatial reports at the same time as independent evidence. With $W_u=\sum_{k:u_k=u}w_k$, the support measure is
\begin{equation}
n_{\mathrm{eff}}(P)=\frac{(\sum_u W_u)^2}{\sum_u W_u^2}.
\label{eq:ess}
\end{equation}
This is a heuristic effective-support statistic, not a claim that dependent samples become independent.

\subsection{Archive mixtures and stale support}
The local pool contains sensor $i$ and up to two graph neighbours. Let $g$ be the nonnegative gap between the current bin and the latest local target, using archive targets as fallback when no live record is available. With reference support $n_0=50$ and decay $d=288$ bins, the real-data local implementation uses
\begin{equation}
\eta=\begin{cases}\min(1,n_{\mathrm{eff}}(P)/n_0)\exp(-g/d),&P\ne\varnothing,\\0,&P=\varnothing.\end{cases}
\qquad F=\eta F_P+(1-\eta)F_A,
\label{eq:blend}
\end{equation}
where $F_A$ is the uniformly weighted frozen archive distribution for the local group, with archive fallback. The local query retains at most $256|G_i|$ records in archive order, where $G_i$ is the sensor group; this is a group limit rather than an equal per-sensor allocation.

The uncapped SUMO implementation also includes a global live pool $Q$. Its live distribution is $\lambda F_P+(1-\lambda)F_Q$, where $\lambda=n_{\mathrm{eff}}(P)/(n_{\mathrm{eff}}(P)+n_0)$. Equation~\eqref{eq:blend} then uses global support for $\eta$ and a sensor-specific archive, falling back to the full archive if fewer than 50 sensor records are available. Both paths implement the same feedback-aware principle, but their pooling estimands and sampling rates differ and are reported separately.

\subsection{Signed-tail interval construction}
Let $z$ denote issue-time missingness, recovered from the stored availability context. Bounded inflation is
\begin{equation}
\kappa=\min\{3,1+0.05\min(g/d,2)+0.05z\}.
\end{equation}
For empirical weighted quantile $Q_p(F)=\inf\{q:F(q)\ge p\}$ and nominal miscoverage $\alpha$, the locked fully asymmetric interval is
\begin{align}
L_{i,t,h}&=\max\{0,m_{i,t,h}+s_{i,t,h}\kappa Q_{\alpha/2}(F)\},\\
U_{i,t,h}&=\max\{L_{i,t,h},m_{i,t,h}+s_{i,t,h}\kappa Q_{1-\alpha/2}(F)\}.
\label{eq:interval}
\end{align}
The candidate uses asymmetric blend one, so no symmetric rolling-radius component remains. Multiplication of signed tails by $\kappa$ need not strictly enlarge both sides when both tail quantiles have the same sign; the equations specify the implemented transformation rather than asserting monotone widening in every case.


\subsection{Predictor input, training objective and noncrossing head}
Let $o^x_{i,u}(t)$ indicate that the value at timestamp $u$ is valid and released to the input store by issue time $t$. Its age feature is the time since the most recent available observation, capped at one day and divided by 288. Together with training-standardised past-filled speed and four calendar channels, this produces $X_t\in\mathbb R^{12\times |V|\times7}$. The validity channel records actual availability even when filling provides a numerical speed. Thus numerical completeness of $X_t$ does not imply that the detector is currently reporting.

For FAR-GW, the learned graph support is
\begin{equation}
A_{\rm adaptive}=\operatorname{softmax}_{\rm row}\{\operatorname{ReLU}(E_1E_2)\},
\end{equation}
where $E_1,E_2$ are trainable node factors. Gated temporal features have the form $\tanh(W_f*X)\odot\sigma(W_g*X)$, with dilated temporal convolutions and graph diffusion projections. These backbone operations generate three raw head outputs $a,b,c$, converted to
\begin{equation}
\widehat q_{.05}=b-\operatorname{softplus}(a),\quad
\widehat q_{.50}=b,\quad
\widehat q_{.95}=b+\operatorname{softplus}(c).
\label{eq:head}
\end{equation}
Equation~\eqref{eq:head} prevents quantile crossing by construction; it does not guarantee empirical coverage. Other backbone adapters preserve the same noncrossing output interface.

With valid training target index set $\mathcal M$ and $\mathcal P=\{.05,.50,.95\}$, training minimises masked pinball loss
\begin{equation}
\mathcal J(\theta)=\frac{1}{3|\mathcal M|}\sum_{p\in\mathcal P}\sum_{(i,t,h)\in\mathcal M}
\rho_p(y_{i,t+h}-\widehat q_{p,i,t,h}),\quad
\rho_p(e)=\max\{pe,(p-1)e\}.
\end{equation}
The implementation evaluates this loss in standardised training units. Empty-validity batches are skipped, and tuning loss selects the checkpoint without test outcomes. The fitted head supplies the centre and scale in Equation~\eqref{eq:scale}; received residuals, rather than further predictor training, update the uncertainty layer during replay.

\subsection{Identifiability and computational assumptions}
The empirical mixture requires a nonempty valid initial archive. Its normalised weights are computed after subtracting the largest log weight for numerical stability, which leaves relative weights unchanged. Zero-mass entries are excluded before quantile lookup. Timestamp-block support reduces the tendency to treat simultaneous neighbouring records as separate independent observations, but is not an independence proof. No finite-sample correction beyond the stated empirical inverse-distribution quantile is applied.

Permanent feedback loss removes information about realised errors. If loss depends on unseen speed, the received residual distribution need not equal the full target distribution. Neither spatial pooling nor historical fallback identifies that difference without additional assumptions. The MNAR stress therefore measures sensitivity to the specified low-speed suppression process. Equations~\eqref{eq:blend} and \eqref{eq:interval} describe the executable heuristic; they do not establish arbitrary-missingness coverage or recover outcomes that the service never received.
